% Options for packages loaded elsewhere
\PassOptionsToPackage{unicode}{hyperref}
\PassOptionsToPackage{hyphens}{url}
\documentclass[
]{article}
\usepackage{xcolor}
\usepackage[margin=1in]{geometry}
\usepackage{amsmath,amssymb}
\setcounter{secnumdepth}{-\maxdimen} % remove section numbering
\usepackage{iftex}
\ifPDFTeX
  \usepackage[T1]{fontenc}
  \usepackage[utf8]{inputenc}
  \usepackage{textcomp} % provide euro and other symbols
\else % if luatex or xetex
  \usepackage{unicode-math} % this also loads fontspec
  \defaultfontfeatures{Scale=MatchLowercase}
  \defaultfontfeatures[\rmfamily]{Ligatures=TeX,Scale=1}
\fi
\usepackage{lmodern}
\ifPDFTeX\else
  % xetex/luatex font selection
\fi
% Use upquote if available, for straight quotes in verbatim environments
\IfFileExists{upquote.sty}{\usepackage{upquote}}{}
\IfFileExists{microtype.sty}{% use microtype if available
  \usepackage[]{microtype}
  \UseMicrotypeSet[protrusion]{basicmath} % disable protrusion for tt fonts
}{}
\makeatletter
\@ifundefined{KOMAClassName}{% if non-KOMA class
  \IfFileExists{parskip.sty}{%
    \usepackage{parskip}
  }{% else
    \setlength{\parindent}{0pt}
    \setlength{\parskip}{6pt plus 2pt minus 1pt}}
}{% if KOMA class
  \KOMAoptions{parskip=half}}
\makeatother
\usepackage{color}
\usepackage{fancyvrb}
\newcommand{\VerbBar}{|}
\newcommand{\VERB}{\Verb[commandchars=\\\{\}]}
\DefineVerbatimEnvironment{Highlighting}{Verbatim}{commandchars=\\\{\}}
% Add ',fontsize=\small' for more characters per line
\usepackage{framed}
\definecolor{shadecolor}{RGB}{248,248,248}
\newenvironment{Shaded}{\begin{snugshade}}{\end{snugshade}}
\newcommand{\AlertTok}[1]{\textcolor[rgb]{0.94,0.16,0.16}{#1}}
\newcommand{\AnnotationTok}[1]{\textcolor[rgb]{0.56,0.35,0.01}{\textbf{\textit{#1}}}}
\newcommand{\AttributeTok}[1]{\textcolor[rgb]{0.13,0.29,0.53}{#1}}
\newcommand{\BaseNTok}[1]{\textcolor[rgb]{0.00,0.00,0.81}{#1}}
\newcommand{\BuiltInTok}[1]{#1}
\newcommand{\CharTok}[1]{\textcolor[rgb]{0.31,0.60,0.02}{#1}}
\newcommand{\CommentTok}[1]{\textcolor[rgb]{0.56,0.35,0.01}{\textit{#1}}}
\newcommand{\CommentVarTok}[1]{\textcolor[rgb]{0.56,0.35,0.01}{\textbf{\textit{#1}}}}
\newcommand{\ConstantTok}[1]{\textcolor[rgb]{0.56,0.35,0.01}{#1}}
\newcommand{\ControlFlowTok}[1]{\textcolor[rgb]{0.13,0.29,0.53}{\textbf{#1}}}
\newcommand{\DataTypeTok}[1]{\textcolor[rgb]{0.13,0.29,0.53}{#1}}
\newcommand{\DecValTok}[1]{\textcolor[rgb]{0.00,0.00,0.81}{#1}}
\newcommand{\DocumentationTok}[1]{\textcolor[rgb]{0.56,0.35,0.01}{\textbf{\textit{#1}}}}
\newcommand{\ErrorTok}[1]{\textcolor[rgb]{0.64,0.00,0.00}{\textbf{#1}}}
\newcommand{\ExtensionTok}[1]{#1}
\newcommand{\FloatTok}[1]{\textcolor[rgb]{0.00,0.00,0.81}{#1}}
\newcommand{\FunctionTok}[1]{\textcolor[rgb]{0.13,0.29,0.53}{\textbf{#1}}}
\newcommand{\ImportTok}[1]{#1}
\newcommand{\InformationTok}[1]{\textcolor[rgb]{0.56,0.35,0.01}{\textbf{\textit{#1}}}}
\newcommand{\KeywordTok}[1]{\textcolor[rgb]{0.13,0.29,0.53}{\textbf{#1}}}
\newcommand{\NormalTok}[1]{#1}
\newcommand{\OperatorTok}[1]{\textcolor[rgb]{0.81,0.36,0.00}{\textbf{#1}}}
\newcommand{\OtherTok}[1]{\textcolor[rgb]{0.56,0.35,0.01}{#1}}
\newcommand{\PreprocessorTok}[1]{\textcolor[rgb]{0.56,0.35,0.01}{\textit{#1}}}
\newcommand{\RegionMarkerTok}[1]{#1}
\newcommand{\SpecialCharTok}[1]{\textcolor[rgb]{0.81,0.36,0.00}{\textbf{#1}}}
\newcommand{\SpecialStringTok}[1]{\textcolor[rgb]{0.31,0.60,0.02}{#1}}
\newcommand{\StringTok}[1]{\textcolor[rgb]{0.31,0.60,0.02}{#1}}
\newcommand{\VariableTok}[1]{\textcolor[rgb]{0.00,0.00,0.00}{#1}}
\newcommand{\VerbatimStringTok}[1]{\textcolor[rgb]{0.31,0.60,0.02}{#1}}
\newcommand{\WarningTok}[1]{\textcolor[rgb]{0.56,0.35,0.01}{\textbf{\textit{#1}}}}
\usepackage{graphicx}
\makeatletter
\newsavebox\pandoc@box
\newcommand*\pandocbounded[1]{% scales image to fit in text height/width
  \sbox\pandoc@box{#1}%
  \Gscale@div\@tempa{\textheight}{\dimexpr\ht\pandoc@box+\dp\pandoc@box\relax}%
  \Gscale@div\@tempb{\linewidth}{\wd\pandoc@box}%
  \ifdim\@tempb\p@<\@tempa\p@\let\@tempa\@tempb\fi% select the smaller of both
  \ifdim\@tempa\p@<\p@\scalebox{\@tempa}{\usebox\pandoc@box}%
  \else\usebox{\pandoc@box}%
  \fi%
}
% Set default figure placement to htbp
\def\fps@figure{htbp}
\makeatother
\setlength{\emergencystretch}{3em} % prevent overfull lines
\providecommand{\tightlist}{%
  \setlength{\itemsep}{0pt}\setlength{\parskip}{0pt}}
\usepackage{bookmark}
\IfFileExists{xurl.sty}{\usepackage{xurl}}{} % add URL line breaks if available
\urlstyle{same}
% fallback for those not using the hyperref driver hyperxmp:
\makeatletter
\@ifundefined{xmpquote}{\newcommand{\xmpquote}[1]{#1}}{}
\makeatother
\hypersetup{
  pdftitle={t-test-function.R},
  pdfauthor={tylerwiederich},
  hidelinks,
  pdfcreator={LaTeX via pandoc}}

\title{t-test-function.R}
\author{tylerwiederich}
\date{2026-11-06}

\begin{document}
\maketitle

\begin{Shaded}
\begin{Highlighting}[]
\FunctionTok{set.seed}\NormalTok{(}\DecValTok{383}\NormalTok{)}
\NormalTok{y }\OtherTok{\textless{}{-}} \FunctionTok{rnorm}\NormalTok{(}\AttributeTok{n=}\DecValTok{30}\NormalTok{, }\AttributeTok{mean=}\DecValTok{20}\NormalTok{, }\AttributeTok{sd=}\DecValTok{3}\NormalTok{)}

\CommentTok{\# H0: µ = 20 vs. HA: π ≠ 20}

\CommentTok{\# t{-}statistic}
\NormalTok{xbar }\OtherTok{\textless{}{-}} \FunctionTok{mean}\NormalTok{(y)}
\NormalTok{s2 }\OtherTok{\textless{}{-}} \FunctionTok{var}\NormalTok{(y)}
\NormalTok{n }\OtherTok{\textless{}{-}} \FunctionTok{length}\NormalTok{(y)}

\NormalTok{t.stat }\OtherTok{\textless{}{-}}\NormalTok{ (xbar}\DecValTok{{-}20}\NormalTok{) }\SpecialCharTok{/} \FunctionTok{sqrt}\NormalTok{(s2}\SpecialCharTok{/}\NormalTok{n)}
\NormalTok{p.val }\OtherTok{\textless{}{-}} \DecValTok{2}\SpecialCharTok{*}\FunctionTok{pt}\NormalTok{(}\FunctionTok{abs}\NormalTok{(t.stat), }\AttributeTok{df=}\NormalTok{n}\DecValTok{{-}1}\NormalTok{, }\AttributeTok{lower.tail =}\NormalTok{ F)}
\NormalTok{ci }\OtherTok{\textless{}{-}}\NormalTok{ xbar }\SpecialCharTok{+} \FunctionTok{c}\NormalTok{(}\SpecialCharTok{{-}}\DecValTok{1}\NormalTok{,}\DecValTok{1}\NormalTok{)}\SpecialCharTok{*}\FunctionTok{qt}\NormalTok{(}\AttributeTok{p=}\DecValTok{1}\FloatTok{{-}0.05}\SpecialCharTok{/}\DecValTok{2}\NormalTok{, }\AttributeTok{df=}\NormalTok{n}\DecValTok{{-}1}\NormalTok{)}\SpecialCharTok{*}\FunctionTok{sqrt}\NormalTok{(s2}\SpecialCharTok{/}\NormalTok{n)}

\CommentTok{\# Print results}
\FunctionTok{cbind}\NormalTok{(xbar, t.stat, p.val, }\AttributeTok{lcl=}\NormalTok{ci[}\DecValTok{1}\NormalTok{], }\AttributeTok{ucl=}\NormalTok{ci[}\DecValTok{2}\NormalTok{])}
\end{Highlighting}
\end{Shaded}

\begin{verbatim}
##          xbar     t.stat    p.val      lcl      ucl
## [1,] 19.34494 -0.8588002 0.397494 17.78492 20.90496
\end{verbatim}

\begin{Shaded}
\begin{Highlighting}[]
\CommentTok{\# t.test function}
\FunctionTok{t.test}\NormalTok{(y, }\AttributeTok{mu =} \DecValTok{20}\NormalTok{, }\AttributeTok{alternative =} \StringTok{\textquotesingle{}two.sided\textquotesingle{}}\NormalTok{)}
\end{Highlighting}
\end{Shaded}

\begin{verbatim}
## 
##  One Sample t-test
## 
## data:  y
## t = -0.8588, df = 29, p-value = 0.3975
## alternative hypothesis: true mean is not equal to 20
## 95 percent confidence interval:
##  17.78492 20.90496
## sample estimates:
## mean of x 
##  19.34494
\end{verbatim}

\begin{Shaded}
\begin{Highlighting}[]
\CommentTok{\# How close is our rejection rate to the significance level?}
\NormalTok{nsim }\OtherTok{\textless{}{-}} \DecValTok{10000}
\NormalTok{reject }\OtherTok{\textless{}{-}} \FunctionTok{numeric}\NormalTok{(nsim)}
\ControlFlowTok{for}\NormalTok{(i }\ControlFlowTok{in} \DecValTok{1}\SpecialCharTok{:}\NormalTok{nsim) \{}
\NormalTok{  y }\OtherTok{\textless{}{-}} \FunctionTok{rnorm}\NormalTok{(}\AttributeTok{n=}\DecValTok{30}\NormalTok{, }\AttributeTok{mean=}\DecValTok{20}\NormalTok{, }\AttributeTok{sd=}\DecValTok{3}\NormalTok{)}
\NormalTok{  test }\OtherTok{\textless{}{-}} \FunctionTok{t.test}\NormalTok{(y, }\AttributeTok{mu=}\DecValTok{20}\NormalTok{, }\AttributeTok{alternative =} \StringTok{\textquotesingle{}two.sided\textquotesingle{}}\NormalTok{)}
\NormalTok{  p.val }\OtherTok{\textless{}{-}}\NormalTok{ test}\SpecialCharTok{$}\NormalTok{p.value}
\NormalTok{  reject[i] }\OtherTok{\textless{}{-}} \FunctionTok{ifelse}\NormalTok{(p.val }\SpecialCharTok{\textless{}} \FloatTok{0.05}\NormalTok{, }\DecValTok{1}\NormalTok{, }\DecValTok{0}\NormalTok{)}
\NormalTok{\}}

\FunctionTok{mean}\NormalTok{(reject)}
\end{Highlighting}
\end{Shaded}

\begin{verbatim}
## [1] 0.0523
\end{verbatim}

\begin{Shaded}
\begin{Highlighting}[]
\CommentTok{\# Two{-}sample t.test}
\CommentTok{\# H0: µ1 = µ2 vs. HA: µ1 ≠ mu2}
\NormalTok{y1 }\OtherTok{\textless{}{-}} \FunctionTok{rnorm}\NormalTok{(}\AttributeTok{n=}\DecValTok{30}\NormalTok{, }\AttributeTok{mean=}\DecValTok{20}\NormalTok{, }\AttributeTok{sd=}\DecValTok{3}\NormalTok{)}
\NormalTok{y2 }\OtherTok{\textless{}{-}} \FunctionTok{rnorm}\NormalTok{(}\AttributeTok{n=}\DecValTok{30}\NormalTok{, }\AttributeTok{mean=}\DecValTok{20}\NormalTok{, }\AttributeTok{sd=}\DecValTok{9}\NormalTok{)}

\FunctionTok{qqnorm}\NormalTok{(y1); }\FunctionTok{qqline}\NormalTok{(y1)}
\end{Highlighting}
\end{Shaded}

\pandocbounded{\includegraphics[keepaspectratio]{t-test-function_files/t-test-function_files/figure-latex/unnamed-chunk-1-1.pdf}}

\begin{Shaded}
\begin{Highlighting}[]
\FunctionTok{qqnorm}\NormalTok{(y2); }\FunctionTok{qqline}\NormalTok{(y2)}
\end{Highlighting}
\end{Shaded}

\pandocbounded{\includegraphics[keepaspectratio]{t-test-function_files/t-test-function_files/figure-latex/unnamed-chunk-1-2.pdf}}

\begin{Shaded}
\begin{Highlighting}[]
\FunctionTok{sd}\NormalTok{(y1); }\FunctionTok{sd}\NormalTok{(y2)}
\end{Highlighting}
\end{Shaded}

\begin{verbatim}
## [1] 2.577274
\end{verbatim}

\begin{verbatim}
## [1] 8.209597
\end{verbatim}

\begin{Shaded}
\begin{Highlighting}[]
\NormalTok{f.stat }\OtherTok{\textless{}{-}} \FunctionTok{var}\NormalTok{(y1)}\SpecialCharTok{/}\FunctionTok{var}\NormalTok{(y2)}
\NormalTok{p.val }\OtherTok{\textless{}{-}} \DecValTok{2}\SpecialCharTok{*}\FunctionTok{min}\NormalTok{(}\FunctionTok{pf}\NormalTok{(}\AttributeTok{q=}\NormalTok{f.stat, }\AttributeTok{df1=}\FunctionTok{length}\NormalTok{(y1)}\SpecialCharTok{{-}}\DecValTok{1}\NormalTok{, }\AttributeTok{df2=}\FunctionTok{length}\NormalTok{(y2)}\SpecialCharTok{{-}}\DecValTok{1}\NormalTok{),}
               \FunctionTok{pf}\NormalTok{(}\AttributeTok{q=}\NormalTok{f.stat, }\AttributeTok{df1=}\FunctionTok{length}\NormalTok{(y1)}\SpecialCharTok{{-}}\DecValTok{1}\NormalTok{, }\AttributeTok{df2=}\FunctionTok{length}\NormalTok{(y2)}\SpecialCharTok{{-}}\DecValTok{1}\NormalTok{, }\AttributeTok{lower.tail =}\NormalTok{ F))}
\NormalTok{ci }\OtherTok{\textless{}{-}}\NormalTok{ f.stat}\SpecialCharTok{/}\FunctionTok{qf}\NormalTok{(}\FunctionTok{c}\NormalTok{(}\FloatTok{0.975}\NormalTok{, }\FloatTok{0.025}\NormalTok{), }\AttributeTok{df1=}\FunctionTok{length}\NormalTok{(y1)}\SpecialCharTok{{-}}\DecValTok{1}\NormalTok{, }\AttributeTok{df2=}\FunctionTok{length}\NormalTok{(y2)}\SpecialCharTok{{-}}\DecValTok{1}\NormalTok{)}
\FunctionTok{cbind}\NormalTok{(f.stat, p.val, }\AttributeTok{lcl=}\NormalTok{ci[}\DecValTok{1}\NormalTok{], }\AttributeTok{ucl=}\NormalTok{ci[}\DecValTok{2}\NormalTok{])}
\end{Highlighting}
\end{Shaded}

\begin{verbatim}
##          f.stat        p.val       lcl       ucl
## [1,] 0.09855477 1.587381e-08 0.0469086 0.2070632
\end{verbatim}

\begin{Shaded}
\begin{Highlighting}[]
\FunctionTok{var.test}\NormalTok{(y1, y2, }\AttributeTok{alternative =} \StringTok{\textquotesingle{}two.sided\textquotesingle{}}\NormalTok{)}
\end{Highlighting}
\end{Shaded}

\begin{verbatim}
## 
##  F test to compare two variances
## 
## data:  y1 and y2
## F = 0.098555, num df = 29, denom df = 29, p-value = 1.587e-08
## alternative hypothesis: true ratio of variances is not equal to 1
## 95 percent confidence interval:
##  0.0469086 0.2070632
## sample estimates:
## ratio of variances 
##         0.09855477
\end{verbatim}

\begin{Shaded}
\begin{Highlighting}[]
\FunctionTok{t.test}\NormalTok{(y1, y2, }\AttributeTok{var.equal =}\NormalTok{ F)}
\end{Highlighting}
\end{Shaded}

\begin{verbatim}
## 
##  Welch Two Sample t-test
## 
## data:  y1 and y2
## t = -0.78612, df = 34.661, p-value = 0.4371
## alternative hypothesis: true difference in means is not equal to 0
## 95 percent confidence interval:
##  -4.425366  1.955402
## sample estimates:
## mean of x mean of y 
##  19.70099  20.93597
\end{verbatim}

\begin{Shaded}
\begin{Highlighting}[]
\CommentTok{\# Rejection rate if equal variance assumption is used?}
\NormalTok{nsim }\OtherTok{\textless{}{-}} \DecValTok{10000}
\NormalTok{reject.eq }\OtherTok{\textless{}{-}} \FunctionTok{numeric}\NormalTok{(nsim)}
\NormalTok{reject.neq }\OtherTok{\textless{}{-}} \FunctionTok{numeric}\NormalTok{(nsim)}
\ControlFlowTok{for}\NormalTok{(i }\ControlFlowTok{in} \DecValTok{1}\SpecialCharTok{:}\NormalTok{nsim) \{}
\NormalTok{  y1 }\OtherTok{\textless{}{-}} \FunctionTok{rnorm}\NormalTok{(}\AttributeTok{n=}\DecValTok{30}\NormalTok{, }\AttributeTok{mean=}\DecValTok{20}\NormalTok{, }\AttributeTok{sd=}\DecValTok{3}\NormalTok{)}
\NormalTok{  y2 }\OtherTok{\textless{}{-}} \FunctionTok{rnorm}\NormalTok{(}\AttributeTok{n=}\DecValTok{30}\NormalTok{, }\AttributeTok{mean=}\DecValTok{20}\NormalTok{, }\AttributeTok{sd=}\DecValTok{9}\NormalTok{)}
\NormalTok{  test.eq }\OtherTok{\textless{}{-}} \FunctionTok{t.test}\NormalTok{(y1, y2, }\AttributeTok{var.equal =}\NormalTok{ T, }\AttributeTok{alternative =} \StringTok{\textquotesingle{}two.sided\textquotesingle{}}\NormalTok{)}
\NormalTok{  test.neq }\OtherTok{\textless{}{-}} \FunctionTok{t.test}\NormalTok{(y1, y2, }\AttributeTok{var.equal =}\NormalTok{ F, }\AttributeTok{alternative =} \StringTok{\textquotesingle{}two.sided\textquotesingle{}}\NormalTok{)}
\NormalTok{  p.val.eq }\OtherTok{\textless{}{-}}\NormalTok{ test.eq}\SpecialCharTok{$}\NormalTok{p.value}
\NormalTok{  p.val.neq }\OtherTok{\textless{}{-}}\NormalTok{ test.neq}\SpecialCharTok{$}\NormalTok{p.value}
  
\NormalTok{  reject.eq[i] }\OtherTok{\textless{}{-}} \FunctionTok{ifelse}\NormalTok{(p.val.eq }\SpecialCharTok{\textless{}} \FloatTok{0.05}\NormalTok{, }\DecValTok{1}\NormalTok{, }\DecValTok{0}\NormalTok{)}
\NormalTok{  reject.neq[i] }\OtherTok{\textless{}{-}} \FunctionTok{ifelse}\NormalTok{(p.val.neq }\SpecialCharTok{\textless{}} \FloatTok{0.05}\NormalTok{, }\DecValTok{1}\NormalTok{, }\DecValTok{0}\NormalTok{)}
\NormalTok{\}}
\FunctionTok{mean}\NormalTok{(reject.eq)}
\end{Highlighting}
\end{Shaded}

\begin{verbatim}
## [1] 0.052
\end{verbatim}

\begin{Shaded}
\begin{Highlighting}[]
\FunctionTok{mean}\NormalTok{(reject.neq)}
\end{Highlighting}
\end{Shaded}

\begin{verbatim}
## [1] 0.0489
\end{verbatim}

\begin{Shaded}
\begin{Highlighting}[]
\FunctionTok{plot}\NormalTok{(}\DecValTok{1}\SpecialCharTok{:}\NormalTok{nsim, }\FunctionTok{cumsum}\NormalTok{(reject.eq)}\SpecialCharTok{/}\DecValTok{1}\SpecialCharTok{:}\NormalTok{nsim, }\AttributeTok{type=}\StringTok{\textquotesingle{}l\textquotesingle{}}\NormalTok{, }
     \AttributeTok{main =} \StringTok{\textquotesingle{}Rejection Rate of True H0\textquotesingle{}}\NormalTok{,}
     \AttributeTok{xlab =} \StringTok{\textquotesingle{}Simulation Number\textquotesingle{}}\NormalTok{,}
     \AttributeTok{ylab =} \StringTok{\textquotesingle{}Cumulative mean rejection rate\textquotesingle{}}\NormalTok{)}
\FunctionTok{lines}\NormalTok{(}\DecValTok{1}\SpecialCharTok{:}\NormalTok{nsim, }\FunctionTok{cumsum}\NormalTok{(reject.neq)}\SpecialCharTok{/}\DecValTok{1}\SpecialCharTok{:}\NormalTok{nsim, }\AttributeTok{type=}\StringTok{\textquotesingle{}l\textquotesingle{}}\NormalTok{, }\AttributeTok{col=}\StringTok{\textquotesingle{}red\textquotesingle{}}\NormalTok{)}
\end{Highlighting}
\end{Shaded}

\pandocbounded{\includegraphics[keepaspectratio]{t-test-function_files/t-test-function_files/figure-latex/unnamed-chunk-1-3.pdf}}

\end{document}
